\label{sec:system-design-body}

This section describes the committed-data pipeline, verification statements,
and SP1 backend. The main design rule is to separate the relation that is
proved in SP1 from the larger Python-only extension checks.

\begin{figure*}[!t]
\centering
\begin{tikzpicture}[
    box/.style={draw, rounded corners=3pt, minimum width=2.7cm, minimum height=0.75cm, align=center, font=\small},
    data/.style={box, fill=blue!8},
    commit/.style={box, fill=green!10},
    proof/.style={box, fill=red!8},
    ext/.style={box, fill=orange!12},
    arr/.style={-{Stealth[length=5pt]}, thick},
    lab/.style={font=\scriptsize, midway, above}
]
\node[data] (dataset) {Offline transition\\dataset $\dataset$};
\node[commit, right=0.7cm of dataset] (merkle) {Merkle\\commitment};
\node[proof, right=0.7cm of merkle] (artifact) {TD / minibatch\\artifact};
\node[proof, right=0.7cm of artifact] (sp1) {SP1 host/guest\\proof backend};
\node[proof, below=0.8cm of sp1] (verify) {Verified proof\\or rejection};
\node[ext, below=0.8cm of artifact] (python) {Python oracle and\\pre-ZK extensions};
\draw[arr] (dataset) -- (merkle);
\draw[arr] (merkle) -- node[lab] {root} (artifact);
\draw[arr] (artifact) -- (sp1);
\draw[arr] (sp1) -- (verify);
\draw[arr] (artifact) -- (python);
\end{tikzpicture}
\caption{Locked Week 5 system architecture. The cryptographic backend proves TD and minibatch-TD relations in SP1. Python verifiers remain as semantic oracles and as pre-ZK extension checks for one-step and short-trace statements.}
\label{fig:system-architecture}
\end{figure*}

\subsection{Threat Model and Non-Goals}

The prover claims that a private transition or minibatch is drawn from the
committed dataset and that the TD outputs are computed correctly. The verifier
knows the dataset root, fixed-point constants, schema version, and claimed
target/loss or batch loss. The prover supplies the private transition contents,
Merkle paths, and TD witness values.

The current backend protects against a prover who tampers with the transition,
Merkle path, leaf index, fixed-point rounding, done-branch semantics,
target-value witness, TD error, claimed target, claimed loss, batch size, or
batch aggregation. It does not protect against dishonest data collection before
commitment, and it does not prove that a full training run produced a final
checkpoint.

\subsection{Canonical Serialization and Commitment}

Each transition is encoded using scale $S=1000$:
\[
\mathrm{fp}(x)=\mathrm{round}(Sx).
\]
The serialized leaf contains the fixed-point state, action, reward,
fixed-point next state, and done flag. The current canonical leaf hash is
SHA-256 over the comma-separated signed integer fields. The Merkle root is
recomputed by iterating over the authentication path and hashing ordered sibling
pairs. The SP1 relation rejects a witness if the serialized transition, leaf
hash, authentication path, or public root are inconsistent.

\subsection{Fixed-Point TD Relation}

The backend uses integer arithmetic:
\[
\mathrm{mul}_{S}(a,b)=\left\lfloor \frac{ab}{S} \right\rfloor,
\qquad
\gamma_{\mathrm{fp}} = 990 .
\]
For a non-terminal transition,
\[
y_{\mathrm{fp}} =
r_{\mathrm{fp}} +
\left\lfloor \frac{\gamma_{\mathrm{fp}} q^{\mathrm{target}}_{\mathrm{fp}}}{S}
\right\rfloor ,
\]
and for a terminal transition $y_{\mathrm{fp}}=r_{\mathrm{fp}}$. The TD error is
\[
\delta_{\mathrm{fp}} =
q^{\mathrm{online}}_{\mathrm{fp}} - y_{\mathrm{fp}}.
\]
The fixed-point SmoothL1 loss is
\[
\ell_{\mathrm{fp}}(\delta)=
\begin{cases}
\left\lfloor \frac{|\delta|^2}{2S} \right\rfloor & \text{if } |\delta| < S,\\
|\delta|-\left\lfloor \frac{S}{2}\right\rfloor & \text{otherwise.}
\end{cases}
\]
The single-transition public claim is accepted only if the recomputed
$y_{\mathrm{fp}}$ and $\ell_{\mathrm{fp}}$ match the public claimed target and
loss.

\subsection{Minibatch TD Relation}

For a minibatch of $k$ items, the SP1 guest repeats the membership and TD checks
for each item. It then computes
\[
L_{\mathrm{batch,fp}} =
\left\lfloor \frac{\sum_{j=1}^{k} \ell_{j,\mathrm{fp}}}{k} \right\rfloor
\]
and checks equality with the public claimed batch loss. The public batch size
must match the number of private witness items.

\subsection{SP1 Backend Split}

The backend is implemented under \texttt{zk\_backend/td\_mvp/sp1/}. The
\texttt{shared} crate defines typed public inputs, private witnesses, Merkle
steps, TD witnesses, fixed-point helpers, hashing, and relation logic. The
\texttt{guest} reads a typed input, calls the shared verifier, and commits the
public output. The \texttt{host} loads JSON fixtures, applies named tamper cases,
runs optional host prechecks, executes the guest, generates proofs for accepted
cases, verifies proofs, and records proving time, verification time, proof size,
and cycle count.

\subsection{Python Oracle and Pre-ZK Extensions}

The Python verifier has two roles. First, it is the semantic oracle used to
cross-check SP1 on valid and tampered TD fixtures. Second, it contains extension
statements for forward TD consistency, one-step SGD updates, and short training
traces. These extensions check checkpoint anchoring, gradient recomputation,
parameter deltas, target-network invariance, checkpoint chaining, target-network
synchronization, and deterministic contiguous sampling rules. They are included
to show the path toward update and trace proofs, but the current ZK claim is
limited to the SP1 TD/minibatch-TD backend.
